\subsection{The binomial construction}
\label{sec:binomial-depth-two-all}

The coefficient Bockstein gives an orientation on a stable module over
$\F_p[t]/t^n$.  Its length modulo $n$ supplies the additive invariant
for every truncation exponent, including when the generic coefficient
algebra is a totally ramified field.  The construction uses all finite
coefficient test modules, so its additivity is established on the full
lattice category.

\begin{theorem}\label{thm:binomial-depth-two-all-exponents}
Let $p$ be a prime, $n\ge2$, $O=\Z_p$ or $\Z_{(p)}$, and
$d\in p^2O$.  Put $A=O[t]/(t^n-d)$.  The kernel
\[
 \ker\!\left(K_0(\Lat_A(S_p))
       \longrightarrow K_0(\Lat_{A[1/p]}(S_p))\right)
\]
contains a subgroup isomorphic to
$\bigoplus_{i\ge1}\Z/n$ that is split in the whole lattice
Grothendieck group.  It has a coordinate retraction annihilating
the degree-$p$ Young permutation span.  Consequently $\Phi_A$
is not surjective in degree $p$.
The same split torsion assertion holds for the cyclic precursors.
\end{theorem}

The case $d=0$ is Theorem~\ref{thm:all-truncations}.
At $p=2$, the cases $d\in8O$ follow from
Corollary~\ref{cor:deep-monogenic-orders}.
We prove the remaining assertions over $O=\Z_p$, and then pass
to $\Z_{(p)}$ by completion.

\subsection{Curvature components and exact residual Tate complexes}

Write $F=t^n-d$, $k=\F_p$, and $R=k[t]/t^n$.
For an $A$-free $C_p$-lattice $E$ with generator matrix $T$,
choose a polynomial lift $\widetilde T$ and put
\[
 \widetilde U=\frac{\widetilde T^p-I}{F},
 \qquad U_E=\widetilde U\bmod F.
\]
The division is possible in $O[t]$.  Cancellation of $F$ shows
that $U_E$ is an $A[C_p]$-endomorphism.
A map is called stably zero if it factors through a finite free
$A[C_p]$-lattice.

\begin{lemma}\label{lem:binomial-all-primary}
The stable class of $U_E$ is independent of the lift and basis
and is natural in $E$.  Stable Hom groups are killed by $p$.
For each irreducible $f\in k[z]$, the integral $f$-primary
summand $E_f$ is $A$-free and is well-defined up to stable
isomorphism.  If $f\ne z$, then $E_f$ and $E_f/pE_f$ have zero
Tate cohomology.
\end{lemma}

\begin{proof}
For an equivariant map $B:E\to E'$, lift it and write
$\widetilde T'\widetilde B-\widetilde B\widetilde T=F\widetilde C$.
Telescoping gives
\[
 U_{E'}B-BU_E
 =\sum_{i=0}^{p-1}T'^{\,p-1-i}CT^i.
\]
This is a group trace and factors through a regular lattice.
The same calculation compares two lifts.  The group trace of
an equivariant map is $p$ times that map.

The usual integral primary projectors for an endomorphism of a
finite free $O$-module are obtained by Chinese remainders modulo
$p^a$ and passage to the inverse limit.  They are $A[C_p]$-linear.
Their summands are $A$-free because $A$ is local, and the
projectors are exact on intertwined short exact sequences.
On the $f$-part some power of $f(U_E)$ is divisible by $p$,
and is therefore stably zero.  Bezout's identity over $k[z]$
makes every stable intertwining map between distinct primary
parts zero.  Comparing two choices using the identity map and
its inverse consequently gives stable isomorphisms of their
corresponding primary parts.

For $f\ne z$, $U_E$ is invertible on $E_f$.  Lift a coefficient
basis adapted to this summand to the local ring
$O[t]_{(p,t)}$.  The off-diagonal action blocks are divisible
by $F$; every term leaving and returning to a diagonal block
of $\widetilde T^p$ is divisible by $F^2$.
The curvature of the diagonal lift is therefore the restricted
curvature, and is a unit in the local ring.

If $(T-I)x=0$ modulo $F$, lift and write
$(\widetilde T-I)x=Fy$.  Multiplication by $N(\widetilde T)$
and cancellation of $F$ give
$\widetilde Ux=N(\widetilde T)y$.
Thus $x$ is a norm boundary.  Reversing the two factors proves
the other Tate exactness assertion.  Multiplication by $p$
on $E_f$ is injective; its coefficient long exact sequence
then proves the same vanishing for $E_f/pE_f$.
\end{proof}

For a compatible action lift on an admissible sequence, choose
upper triangular blocks.  Its curvature is upper triangular
with the endpoint curvatures on the diagonal.  The primary
projectors therefore give admissible sequences of the selected
components.  This observation will supply exactness before any
coordinate is evaluated.

\subsection{A stable module representing all coefficient tests}

Let $E$ be a selected cohomologically trivial component, and put
$\overline E=E/pE$.  Its exact, two-periodic Tate complex has
free $R$-modules as terms and differentials
\[
 D^0=T-I,\qquad D^1=N(T).
\]
Write $(-)^*=\operatorname{Hom}_R(-,R)$ and define
\[
 M_E=\operatorname{coker}
 \bigl((D^0)^*:\overline E^*\longrightarrow\overline E^*\bigr).
\]

The coefficient-of-$t^{n-1}$ pairing identifies $R$ with its
$k$-linear dual, so $R$ is a symmetric Frobenius algebra.
Its finite-module stable category
$\mathcal S_n=\underline{\operatorname{mod}}(R)$ is triangulated,
with triangles induced by short exact sequences and suspension
given by an injective embedding
\citep[Section~13.4, Theorem~13.7, p.~50 of the author's
version]{Buhler2010}.

\begin{lemma}\label{lem:binomial-all-representability}
For every finite $R$-module $X$ there is a natural isomorphism
\begin{equation}\label{eq:binomial-all-representable}
 \widehat H^0(C_p,\overline E\otimes_R X)
 \cong \underline{\operatorname{Hom}}_R(M_E,X).
\end{equation}
The assignment $E\mapsto M_E$ is contravariant.
An admissible sequence of selected components gives an actual
short exact sequence
\begin{equation}\label{eq:binomial-all-transpose-sequence}
 0\longrightarrow M_{E''}\longrightarrow M_E
   \longrightarrow M_{E'}\longrightarrow0,
\end{equation}
and the resulting stable Hom long exact sequence agrees with
the Tate long exact sequence under \eqref{eq:binomial-all-representable}.
\end{lemma}

\begin{proof}
Dualizing the exact Tate complex preserves exactness because
$R$ is self-injective.
Ordinary $\operatorname{Hom}_R(M_E,X)$ identifies with the
kernel of $D^0$ on $\overline E\otimes_R X$.
The dual complex embeds $M_E$ as the image of $(D^1)^*$ in
a free module.  Every map from $M_E$ to a projective extends
across this embedding, since projectives are injective.
Maps factoring through projectives are consequently exactly
the image of $D^1$ on $\overline E\otimes_R X$.  This proves
\eqref{eq:binomial-all-representable}.

An equivariant map $E\to E'$ dualizes to a map of complexes
in the opposite direction and hence gives $M_{E'}\to M_E$.
For an admissible sequence, the dual complexes form a
degreewise split short exact sequence of exact complexes.
The induced map on cycles is surjective: express a quotient
cycle as a differential, lift its preimage, and apply the
differential.  The cokernel sequence is therefore
\eqref{eq:binomial-all-transpose-sequence}.
Applying Hom to these complete resolutions identifies its
connecting map with the Tate connecting map, as asserted.
\end{proof}

\subsection{Lifting every test-module map and its suspension}

Put
\[
 \widetilde R=(O/p^2)[t]/t^n=A/p^2A.
\]
Every finite $R$-module is a direct sum of
$X_i=R/(t^i)$, $1\le i\le n$.  Use the standard lifts
$\widetilde X_i=(O/p^2)[t]/t^i$ and their direct sums.
Every map $X_i\to X_j$ lifts: its value at $1$ is divisible
by $t^{\max(j-i,0)}$, and the remaining coefficients lift
arbitrarily.  This also proves the assertion for matrices
of maps between direct sums.  An isomorphism lifts to an
isomorphism because it is invertible modulo $p$ between
finite free $O/p^2$-modules.

For a test module $X$ these lifts give a coefficient sequence
\[
 0\longrightarrow\overline E\otimes_R X
 \xrightarrow{p}(E/p^2E)\otimes_{\widetilde R}\widetilde X
 \longrightarrow\overline E\otimes_R X\longrightarrow0.
\]
It is exact since $E/p^2E$ is $\widetilde R$-free and
$\widetilde X$ is $O/p^2$-free.  Let $\beta_X$ be its Tate
Bockstein.  Lifting each test-module map proves naturality
on the entire category of finite $R$-modules.  A functorial
choice of lifts of all maps is unnecessary: any lift gives
the required naturality square with the prescribed reduction.
In particular, changes of cyclic decomposition of $X$
do not change the resulting operator.

Set
\[
 Q_p(z)=\frac{N(z)-p}{z-1},\qquad
 h_{\rm even}=I,\qquad h_{\rm odd}=-Q_p(T).
\]
The identity $dh+hd=pI$ makes $h$ an odd Tate operator modulo
$p$.  If a lifted cycle satisfies $dz=pw$, then
$d(hz)=p(z-hw)$.  Division by $p$ proves
\begin{equation}\label{eq:binomial-all-beta-h}
 \beta h+h\beta=I.
\end{equation}
All these operators are natural in both $E$ and $X$.
Define the even operator $P_X=\beta_Xh_X$.

\begin{lemma}\label{lem:binomial-all-boundaries}
The operator $P$ commutes with every Tate connecting map of
an admissible lattice sequence.  It also commutes with the
coefficient connecting maps implementing the stable suspension
in the variable $X$.
\end{lemma}

\begin{proof}
For a lattice sequence, reduce modulo $p^2$, tensor with
$\widetilde X$, and choose an $O/p^2$-linear splitting.
The differentials have upper triangular blocks.  Their
off-diagonal block $Z$ is an odd cochain map modulo $p^2$.
On a cycle with $dz=pw$, division by $p$ therefore gives
$\beta'\partial+\partial\beta''=0$.
The upper block of $dh+hd=pI$ gives
$h'\partial+\partial h''=0$ modulo $p$.
Their product commutes with $\partial$.

For suspension use
\[
 0\longrightarrow X_i\xrightarrow{t^{n-i}}R
 \longrightarrow X_{n-i}\longrightarrow0.
\]
These sequences lift to exact sequences of the standard
$\widetilde R$-modules and are split over $O/p^2$.
The identical two-boundary calculation shows that both
$\beta$ and $h$ anticommute with their coefficient connecting
maps.  Hence $P$ commutes with them.  The middle test module
$R$ has zero Tate cohomology, so these connecting maps
identify the two parities with stable suspension.
Direct sums cover every finite test module.
\end{proof}

By Yoneda and Lemma~\ref{lem:binomial-all-representability},
$P^0_X$ is induced by a unique stable endomorphism
$P_E$ of $M_E$.  The suspension assertion identifies $P^1$
with the shifted endomorphism.  Consequently the first
assertion of Lemma~\ref{lem:binomial-all-boundaries} says that
these endomorphisms commute with all three arrows of the
triangle \eqref{eq:binomial-all-transpose-sequence},
including its connecting arrow.

\subsection{Stable summands and additive length}

The category $\mathcal S_n$ is idempotent complete.
Indeed, every finite $R$-module is a sum of the cyclic
$X_i$.  Remove the free summands $X_n$.
On $X_i$, $i<n$, a map factoring through a free module is
divisible by $t^{n-i}$ and belongs to the radical of its
endomorphism ring.  The projective-factor ideal on a finite
sum is therefore contained in the radical of its finite-dimensional
endomorphism algebra.  Idempotents lift through this nilpotent
ideal and split as module summands.

The function
\[
 \lambda(M)=\dim_k M\pmod n
\]
is well-defined on stable isomorphism classes.
For completeness, stable inverse maps give a split inclusion
$M\to N\oplus Q$ with $Q$ projective; its complementary
idempotent factors through projectives, so its complement
is projective.  Thus a stable isomorphism becomes an actual
isomorphism after adding free $R$-modules.  Such modules
have dimensions divisible by $n$.
The function $\lambda$ is additive on distinguished triangles
because they are induced by short exact sequences.

We also use that a direct summand of a distinguished triangle
is distinguished
\citep[Lemma~1.6 and proof, pp.~821--822]{BalmerSchlichting2001}.
Thus a compatible idempotent on a triangle can be split
without losing the relation needed for $\lambda$.

\subsection{Odd primes}

Suppose $p$ is odd.  The group generator acts trivially on
Tate cohomology, and
\[
 h^2=-Q_p(1)=-p(p-1)/2=0.
\]
Equation~\eqref{eq:binomial-all-beta-h} gives
\[
 P^2=\beta(1-\beta h)h=P-\beta^2h^2=P.
\]
No assertion that $\beta^2=0$ is used.
Let $M_E^+$ denote the image of the stable idempotent $P_E$.
The preceding triangle compatibilities and direct-summand
lemma imply
\[
 \lambda(M_E^+)=\lambda(M_{E'}^+)+\lambda(M_{E''}^+).
\]

For each irreducible $f\in k[z]$, $f\ne z$, define
\begin{equation}\label{eq:binomial-all-odd-coordinate}
 \eta_f(L)=-\lambda(M_{L_f}^+)\in\Z/n.
\end{equation}
These are independent of the curvature choices.
In fact, a map factoring through $A[C_p]^a$ acts as zero on
every test functor: after reduction and tensoring with $X$
the regular module is free over $k[C_p]$.
Stable equivalences of the primary components therefore
give isomorphisms of their representable stable Hom functors,
intertwining the natural Bockstein and $h$.
Yoneda gives compatible stable isomorphisms of their
oriented summands.

For an admissible sequence choose the compatible primary
projectors from Lemma~\ref{lem:binomial-all-primary}.
The full triangle argument now proves additivity of
\eqref{eq:binomial-all-odd-coordinate}.
Thus these are homomorphisms on the full lattice $K_0$.
Each class has finite primary support.  Constant $O$-lattices
have a constant action lift with curvature zero, so every
$\eta_f$ annihilates them.

\subsection{The prime two at valuation two}

Suppose $p=2$ and $d=4b$ with $b\in O^\times$.
Let $E$ be the entire invertible Fitting part of the curvature.
Here $h^2=-I$ and $P$ is not assumed idempotent.
For any standard test module $X_i$, lift to $O[t]/t^i$.
The products of the lifted Tate differentials are
$-4b\widetilde U(t)$, since $i\le n$.
If $D_i z=2w$, then
$D_{i+1}w=-2b\widetilde U(t)z$.
Dividing twice by two gives, on every test module,
\[
 \beta^2=-\overline b\,U(t)=U(t),\qquad
 P^2+P=U(t).
\]
This uses the full curvature polynomial on $X_i$, not merely
its value at $t=0$.

The stable endomorphism $P_E$ of $M_E$ still commutes with
all arrows of the associated triangle.
For an irreducible $q\in\F_2[z]$, $q\ne z,z+1$, let
$M_{E,q}$ be its primary stable summand.
For a triangle, choose a common annihilating polynomial
of the three finite-dimensional stable endomorphisms.
Chinese remainders give one polynomial primary projector
on all three vertices; it is an idempotent endomorphism
of that triangle.  Define
\begin{equation}\label{eq:binomial-all-two-coordinate}
 \eta_q(L)=\lambda(M_{E,q})\in\Z/n.
\end{equation}
The direct-summand lemma proves additivity.
The same stable-equivalence argument used at odd primes
proves choice independence.  These coordinates have finite
support and vanish on constant lattices.
The relation $P(P+I)=U(t)$ places the selected $q$-part
over invertible curvature.

\subsection{Binomial witnesses and their exact annihilation}

Let $O_s/O$ be a finite unramified extension, and choose
units $\alpha,c\in O_s$ with
\begin{equation}\label{eq:binomial-all-parameter}
 \alpha^p+cd=1.
\end{equation}
Put $A_s=O_s\otimes_O A$, let $C_\alpha$ be the companion
of $q_\alpha(z)=(z^p-\alpha^p)/(z-\alpha)$,
and let $a$ be its last coordinate row and $b_0$ its first
coordinate column.  Define
\[
 T_j=
 \begin{pmatrix}
 \alpha&ct^j a\\ t^{n-j}b_0&C_\alpha
 \end{pmatrix},
 \qquad 0\le j\le n.
\]
The companion identity gives
\[
 \det(zI-T_j)=z^p-\alpha^p-ct^n=z^p-1,\qquad
 \widetilde T_j^p-I=c(t^n-d)I.
\]
Thus $L_j=(A_s^p,T_j)$ is a lattice with curvature $c$.

Set $D_-=\operatorname{diag}(1,tI_{p-1})$ and
$D_+=\operatorname{diag}(t,I_{p-1})$.
For $1\le j<n$ they give coefficient-split exact sequences
\[
 0\longrightarrow L_j
 \xrightarrow{(D_-,D_+)}L_{j-1}\oplus L_{j+1}
 \xrightarrow{(D_+,-D_-)}L_j\longrightarrow0.
\]
On each coordinate these are the split sequences
$x\mapsto(x,tx)$, $(v,w)\mapsto tv-w$, or their transposes.
The endpoint lattices are regular: at $j=0$ the unit upper
attachment gives a cyclic companion vector, and at $j=n$
the unit lower attachment does so.
Writing $r_j=[L_j]$, the relations
$2r_j=r_{j-1}+r_{j+1}$ and $r_n=r_0$ yield
\begin{equation}\label{eq:binomial-all-annihilator}
 n\bigl([L_1]-[A_s[C_p]]\bigr)=0.
\end{equation}
This is an actual exact-category relation, independent
of the coordinate calculations.

\subsection{The transpose cokernel and its selected summand}

We give an explicit presentation.
All residue calculations below are over
$R_s=k_s[t]/t^n$, where $k_s=O_s/pO_s$.
Since $\alpha\equiv1\pmod p$, put $m=p-1$ and write
\[
 \overline q_\alpha(z)
 =z^m+q_{m-1}z^{m-1}+\cdots+q_0=(z-1)^m.
\]
Present $M_{L_j}=\operatorname{coker}(T_j-I)^*$ by
generators $x,y_1,\ldots,y_m$.
The top action row gives $\overline c\,t^j y_m=0$.
For $2\le r\le m$, the companion rows give
\[
 y_{r-1}-y_r-q_{r-1}y_m=0.
\]
Eliminating $y_1,\ldots,y_{m-1}$ with these unit-coefficient
relations leaves the first companion-row relation
\[
 t^{n-j}x-\bigl(1+q_{m-1}+\cdots+q_0\bigr)y_m
 =t^{n-j}x-\overline q_\alpha(1)y_m=t^{n-j}x=0.
\]
For $m=1$ the elimination list is empty and the same final
relation holds directly.  Since $\overline c$ is a unit,
\begin{equation}\label{eq:binomial-all-cokernel}
 M_{L_j}\cong R_s/(t^{n-j})\,x\ \oplus\
                    R_s/(t^j)\,y_m.
\end{equation}
The first generator corresponds under
\eqref{eq:binomial-all-representable} at $X=k_s$ to the
even trivial Tate class; the second corresponds to the
even cyclotomic Tate class.

At odd $p$, choose $d=p^v b$, $v\ge2$, and a unit $y\in O_s$.
Set
\[
 \alpha=1+p^{v-1}y,\qquad c=(1-\alpha^p)/d.
\]
The binomial expansion shows that $c$ is a unit and
$\overline c=-\overline y/\overline b$, so every nonzero
residue parameter is available.
On the trivial class $h_{\rm even}=I$ and
$\beta_{\rm odd}=I$, because
\[
 \frac{N(\alpha)}p\equiv1\pmod p.
\]
For the latter identity expand
$N(1+s)=\sum_{\ell=1}^p\binom p\ell s^{\ell-1}$ with
$s\in pO_s$; after division by $p$ every term after the
constant term is divisible by $p$.
On the even cyclotomic class, represented by
$(C_\alpha-I)^{p-2}b_0$ modulo $p$, the map $h_{\rm even}=I$
has odd image in $\operatorname{im}(C_\alpha-I)$, and hence
is zero in cohomology.
Thus $P$ selects the trivial even class and kills the
cyclotomic even class.

Under \eqref{eq:binomial-all-cokernel}, the selected stable
summand is consequently $R_s/(t^{n-j})$.
Indeed, the induced idempotent has diagonal residues $1,0$
in the semisimple quotient of the endomorphism algebra
of these cyclic modules.  Krull--Schmidt and idempotent
lifting determine its image up to stable isomorphism.
For $j=1$ this also covers $n=2$, when the two cyclic
summands have equal length and the selected multiplicity
is one.  If $\overline c$ has degree $s$ and minimal
polynomial $f$, coefficient restriction to $O$ therefore gives
\begin{equation}\label{eq:binomial-all-odd-value}
 \eta_f(L_1)=s\pmod n.
\end{equation}

At $p=2$, $d=4b$, choose an odd prime degree $s>n$ and
a primitive element $\overline y\in k_s^\times$, with
an integral lift $y$.  Put
\[
 \alpha=1+2y,\qquad c=-y(1+y)/b.
\]
Both $y$ and $1+y$ are units and
\eqref{eq:binomial-all-parameter} holds.
On even residual Tate cohomology,
\[
 P_{\rm even}=\operatorname{diag}(1+\overline y,\overline y),
\]
since $h_{\rm even}=I$ and
$\beta_{\rm odd}=(T_j(0)+I)/2$.
The elements $\overline y$ and $1+\overline y$ have
different absolute traces because $s$ is odd, so their
minimal polynomials are distinct.
If $q$ is the minimal polynomial of $\overline y$,
its primary projector in \eqref{eq:binomial-all-cokernel}
selects $R_s/(t^j)y_m$.
The same endomorphism-algebra argument just used makes
this a statement about the actual stable summand.
Thus
\begin{equation}\label{eq:binomial-all-two-value}
 \eta_q(L_1)=s\pmod n.
\end{equation}

\subsection{Induction, retractions, and the two coefficient bases}

\begin{proof}[Proof of Theorem~\ref{thm:binomial-depth-two-all-exponents}]
For odd $p$ and $d\ne0$, choose distinct prime degrees
$s_i>n(p-1)$ and nonzero residue parameters
$\overline c_i$ of those degrees.
Let $f_i$ be their minimal polynomials and restrict the
corresponding $L_1$ from $A_{s_i}$ to $A$, obtaining $L_i$.
The scalar orbits of the parameters under $\F_p^\times$
are free.  If $a\overline c_i$ is Frobenius-conjugate to
$\overline c_i$, taking the field norm gives $a^{s_i}=1$;
coprimality with $p-1$ forces $a=1$.
Different degrees give disjoint orbits.

Changing the cyclic generator to its $a$th power multiplies
the stable curvature by $a$, as follows by expanding
$(I+F\widetilde U)^a$ and dividing by $F$.
Mackey restriction of induction from $C_p$ to $S_p$ supplies
each normalizer twist once.  Every other intersection is
trivial and gives a regular cyclic lattice.
All nonidentity normalizer twists have zero $f_i$-primary
part before applying $P$, so no invariance of $P$ under
generator changes is required.
For
\[
 x_i=[\operatorname{Ind}_{C_p}^{S_p}L_i]
                            -s_i[A[S_p]]
\]
we obtain
\[
 \bigl(\eta_{f_j}\operatorname{Res}_{C_p}^{S_p}\bigr)(x_i)
 =\begin{cases}s_i\pmod n,&i=j,\\0,&i\ne j.\end{cases}
\]
Normalize each coordinate by $s_i^{-1}\in\Z/n$.
Equation~\eqref{eq:binomial-all-annihilator}, preserved by
coefficient restriction and induction, gives $nx_i=0$.
The normalized coordinates, with finite support on each
class, are a retraction of the homomorphism sending the
standard generators of $\bigoplus_i\Z/n$ to the $x_i$.
This proves exact order $n$ and splitting in the whole group.
The cyclic assertion follows before induction.

For $p=2$ and $v_2(d)=2$, choose distinct odd prime degrees
$s_i>n$ and the parameters in
\eqref{eq:binomial-all-two-value}.
The two primary parameters for each witness have degree
$s_i$, so different degrees separate all witnesses.
The values $s_i$ are units modulo $n$.
The same annihilation relation and normalized coordinates
give the split direct sum.  Here $C_2=S_2$.

The remaining nonzero case is $p=2$ and $d=2^v b$ with
$v\ge3$ and $b$ a unit.  Choose the unramified residue parameters
of Theorem~\ref{thm:all-truncations}, and set
$\alpha=1+2^{v-1}y$, $c=(1-\alpha^2)/d$.
Then $c$ is a unit with residue $-\overline y/\overline b$,
and $\alpha\equiv1\pmod4$.  The matrices $T_j$ above reduce
modulo four to the truncated-polynomial witnesses with the same
parameter $c$.  The coefficient comparison of
Theorem~\ref{thm:congruent-coefficients}, at precision three for
$C_2$, transports those witnesses to exact actions with the same
reduction modulo four.  Its rigidity argument identifies the
transported actions with $T_j$, since four has valuation strictly
larger than $v_2(|C_2|)=1$.  Thus the transported coordinates take
the same values on these explicit companions.  The mesh relation
\eqref{eq:binomial-all-annihilator} gives their exact order $n$.
This also supplies explicit generators in the deeper case.

For the constructed witnesses, $t$ is invertible after
inverting $p$, since $d\ne0$.  The diagonal maps relating
successive $L_j$ become isomorphisms, and $L_0$ is regular.
Hence the classes vanish generically.
Every Young permutation lattice, and its restriction to
$C_p$, is extended from $O$.  Its curvature is zero, so
all coordinates annihilate the Young span.
This proves non-surjectivity in degree $p$.

The cases $d=0$ and $p=2,\ d\in8O$ were covered at the
start by the truncation and deep-monogenic results.
Finally let $O=\Z_{(p)}$.  The completion of $A$ is the
corresponding binomial algebra over $\Z_p$.
Theorem~\ref{thm:completion-generic-kernel} lifts the
completed split subgroup through the generic-kernel
isomorphism.  Compose its completed retraction with
completion to obtain a retraction on the whole integral
Grothendieck group.  Young permutation lattices complete
to the corresponding Young lattices, so this retraction
still annihilates their span.

The companion witnesses constructed above can also be chosen
over $\Z_{(p)}$ itself.  For each residue parameter choose a
monic lift of its separable minimal polynomial and let $O_s$
be the resulting finite free local etale $\Z_{(p)}$-algebra.
Its completion is the unramified extension used above.
For odd $p$, and also for $p=2$ with $v_2(d)\ge3$, the formulas
$\alpha=1+p^{v-1}y$ and $c=(1-\alpha^p)/d$ lie in $O_s$:
the numerator is divisible by $p^v$ and $d/p^v$ is a unit.
For $p=2$, $v_2(d)=2$, use
$\alpha=1+2y$ and $c=-y(1+y)/(d/4)$.
The companion matrices, coefficient restriction, induction,
and mesh sequences are then defined over $\Z_{(p)}$.
Their completed detector values are exactly those computed
above, so they give the same sections over the localized integers.
This proves all assertions.
\end{proof}

\begin{example}\label{ex:quintic-depth-two-twentyfive}
Let $p=n=5$, $A=\Z_5[t]/(t^5-25)$, $\alpha=-4$ and $c=41$.
Then $\alpha^5+25c=1$, and
\[
 q_\alpha(z)=z^4-4z^3+16z^2-64z+256.
\]
With $C_\alpha$ its companion, last coordinate row $a$,
and first coordinate column $b_0$, the matrix
\[
 T=\begin{pmatrix}-4&41t\,a\\t^4b_0&C_\alpha\end{pmatrix}
\]
satisfies $T^5=I$.
Its curvature has residue $1$, whose scalar orbit under
$\F_5^\times$ is free.  The degree-one coordinate in
\eqref{eq:binomial-all-odd-value} shows that
$[\operatorname{Ind}_{C_5}^{S_5}L]-[A[S_5]]$ has exact
order five.
If $\pi^5=5$, then $t=\pi^2$ and
$A=\Z_5[\pi^2]$ has index $25$ in the totally ramified
maximal order $\Z_5[\pi]$.
Thus the conclusion includes the exponent equal to
the residue characteristic and a totally ramified generic field.
\end{example}
